% Persian cover letter — Quranic Emotion Spectrum
% Karim Jabbary & Ali Jabbary — University of Urmia
% Target: Pizhuhish'ha-yi Zabanshenakhti-yi Qur'an, University of Isfahan
% Compile: xelatex cover_letter.tex   (uses the project template.tex via pandoc)

\documentclass[12pt,a4paper]{article}

\usepackage[a4paper,top=2.6cm,bottom=2.6cm,left=2.6cm,right=2.6cm]{geometry}

% Fonts (XeLaTeX) ------------------------------------------------------
\usepackage{fontspec}
\defaultfontfeatures{Ligatures=TeX}

\setmainfont{EB Garamond}[
  Path        = ../fonts/,
  Extension   = .ttf,
  UprightFont = ebgaramond-regular,
  Numbers     = Lining,
]

\newfontfamily\arabicfont{Amiri}[
  Path        = ../fonts/,
  Extension   = .ttf,
  UprightFont = amiri-regular,
  BoldFont    = amiri-bold,
  Script      = Arabic,
]
\newcommand{\textarabic}[1]{{\arabicfont #1}}
\DeclareRobustCommand{\arb}[1]{{\arabicfont #1}}

% xepersian — must come AFTER fontspec
\usepackage{xepersian}
\settextfont{Vazirmatn}[
  Path        = ../fonts/,
  Extension   = .ttf,
  UprightFont = vazirmatn-regular,
  BoldFont    = vazirmatn-bold,
  Script      = Arabic,
  Scale       = 1.0,
]
\setlatintextfont{EB Garamond}[
  Path        = ../fonts/,
  Extension   = .ttf,
  UprightFont = ebgaramond-regular,
]
\setdigitfont{Vazirmatn}[
  Path        = ../fonts/,
  Extension   = .ttf,
  UprightFont = vazirmatn-regular,
  Script      = Arabic,
]

\setcounter{secnumdepth}{-1}

\usepackage{setspace}
\setstretch{1.35}
\setlength{\parskip}{0.45\baselineskip}
\setlength{\parindent}{1.4em}

\usepackage{enumitem}
\setlist{leftmargin=1.8em,topsep=0.3ex,itemsep=0.3ex}

\usepackage[hidelinks]{hyperref}
\hypersetup{
  pdftitle={Cover Letter — Quranic Anger Spectrum (Persian)},
  pdfauthor={Karim Jabbary; Ali Jabbary},
  colorlinks=false,
}

% Pandoc compatibility (in case compiled through pandoc template)
\providecommand{\tightlist}{\setlength{\itemsep}{0pt}\setlength{\parskip}{0pt}}

\widowpenalty=10000
\clubpenalty=10000

\pagestyle{empty}

\begin{document}

\begin{flushright}\small
\textbf{تاریخ}: ۲۶ اردیبهشت ۱۴۰۵ (۱۶ مه ۲۰۲۶)
\end{flushright}

\begin{center}\large\bfseries
به نام خدا
\end{center}

\vspace{0.4em}

\noindent\textbf{به}: سردبیرِ محترمِ نشریه‌ی \emph{پژوهش‌های زبان‌شناختیِ قرآن}، دانشگاه اصفهان\\
\textbf{از}: کریم جباری (نویسنده‌ی مسئول) و علی جباری، دانشگاه ارومیه\\
\textbf{موضوع}: ارسالِ مقاله‌ی پژوهشی-تحلیلی جهتِ بررسی و داوری

\vspace{0.4em}\hrule\vspace{0.6em}

\noindent جنابِ آقای دکتر سردبیرِ محترمِ نشریه‌ی پژوهش‌های زبان‌شناختیِ قرآن،

\noindent با سلام و احترام،

به پیوست، مقاله‌ای پژوهشی-تحلیلی با عنوانِ «اوجِ طیفِ خشم در قرآن کریم: قرائتی تطبیقی-تفسیری از مرحله‌های فعّال، متراکم، و پیامدی»، جهتِ بررسی و داوری در آن نشریه‌ی وزین تقدیم می‌گردد. موضوعِ این پژوهش --- تحلیلِ شبکه‌ی واژگانیِ خشم در پرتوِ معناشناسیِ میدانی و نظریه‌ی استعاره‌ی مفهومی --- با گستره‌ی موضوعیِ نشریه‌ی شما در مطالعاتِ زبان‌شناختی و معناشناختیِ قرآن کریم تناسبِ مستقیم دارد.

\textbf{خلاصه‌ی سهمِ پژوهش.} بنده در این پژوهش، با تکیه بر پیکره‌ی ریخت‌شناختیِ قرآنیِ دوکس (۲۰۱۱) و در مجموع ۳۱۲ بسامد از چهارده ریشه‌ی کانونی، پیوستاری شش‌مرحله‌ای از طیفِ خشمِ قرآنی صورت‌بندی کرده و تحلیلِ تفصیلی را بر سه مرحله‌ی اوج --- فعّال، متراکم و پیامدی --- متمرکز ساخته‌ام. در مرحله‌ی پیامدی، کدگذاریِ تفسیر-بنیادِ هر ۱۴۵ بسامد، با کاپای کوهنِ \lr{$\kappa = 0.79$}، نشان می‌دهد که علّیتِ این مرحله یک‌سویه نیست و در کنارِ مسیرِ هیجانی، مسیرِ ساختاریِ مستقلی (طمع، استکبار، عُلوّ فی الأرض) نیز در کار است. هم‌چنین تلاش شده است سه‌گانه‌ی غیظ-کَظم-تَمَیُّز (آل عمران: ۱۳۴ و ۱۱۹؛ مُلک: ۸) به‌مثابهِ عملیاتی‌سازیِ چهار-گامیِ استعاره‌ی \emph{خشم به‌مثابهِ ظرفِ تحتِ فشار} بازخوانی شود.

\textbf{هم‌خوانی با گستره‌ی نشریه.} سه ویژگیِ این مقاله، آن را به‌طورِ خاص با مأموریتِ \emph{پژوهش‌های زبان‌شناختیِ قرآن} هم‌سو می‌کند. نخست، تحلیلِ تطبیقیِ نُه مفسّرِ کانونیِ شیعی و سنّی (المیزان، مجمع‌البیان، الکشّاف، التحریر و التنویر، جامع‌البیان، الجامعِ قرطبی، ابن‌کثیر، فی ظِلال القرآن و المنار)، با مراجعه‌ی موضعی به لطائف‌الإشاراتِ قشیری در بحثِ کَظمِ غیظ. دوم، گفت‌وگوی نظام‌مند با میراثِ معناشناسیِ شناختیِ ایرانی-اسلامی، به‌ویژه آثارِ قائمی‌نیا (۱۳۹۰، ۱۳۹۳) و پاکتچی (۱۳۸۷ به بعد)، و در امتدادِ پژوهش‌های اخیرِ نَریمانی و همکاران (۱۴۰۰) و قاری‌زاده و همکاران (۱۴۰۲). سوم، تکیه بر یک خط‌لوله‌ی پردازشیِ بازتولیدپذیر به زبانِ پایتون که، در حدِّ بررسیِ ادبیات، گامی روش‌شناختی در پیکره‌خوانیِ فارسی‌محورِ متنِ قرآن به‌شمار می‌آید.

\textbf{اظهارنامه‌های اخلاقی.} این مقاله اثرِ اصیلِ نویسندگان است؛ پیش از این در نشریه‌ی دیگری چاپ نشده، هم‌اکنون در دستِ داوریِ نشریه‌ی دیگری نیست و برای ارسالِ هم‌زمان به نشریه‌ی دیگری در نظر گرفته نشده است. هر دو نویسنده در طراحی، نگارش و بازنگری مشارکتِ معنادار داشته‌اند و هیچ‌یک تعارضِ منافع برای اعلام ندارند. مجوزهای پیکره‌ای (\lr{QAC, GPL}؛ تنزیل، \lr{CC-BY-ND}) در بدنه‌ی مقاله تصریح شده است.

\textbf{پیشنهادِ داور} (با احترام به اختیارِ تامِ هیئتِ تحریریه):

\begin{itemize}\setlength{\itemsep}{0pt}
  \item جنابِ آقای دکتر علیرضا قائمی‌نیا (پژوهشگاهِ فرهنگ و اندیشه‌ی اسلامی) --- معناشناسیِ شناختیِ قرآن و استعاره‌ی مفهومی.
  \item جنابِ آقای دکتر احمد پاکتچی (دانشگاه امام صادق علیه‌السلام) --- معناشناسیِ تاریخی-تفسیریِ قرآن.
  \item سرکار خانم دکتر زهرا نَریمانی (دانشگاه تهران) --- معناشناسیِ تفسیریِ واژگانِ هیجانی در قرآن.
\end{itemize}

از توجه و وقتِ ارزشمندِ جنابعالی و هیئتِ تحریریه‌ی محترم، صمیمانه سپاس‌گزارم. این‌جانب آماده‌ام هرگونه اصلاحیه‌ی پیشنهادیِ داورانِ محترم را با گشاده‌رویی بپذیرم.

\vspace{0.6em}
\noindent با تشکر و احترام،

\vspace{0.8em}
\noindent\textbf{کریم جباری} (نویسنده‌ی مسئول، به نیابت از هر دو نویسنده)\\
گروهِ زبان و ادبیاتِ عرب، دانشگاه ارومیه\\
رایانامه: \lr{\texttt{k.jabbary@urmia.ac.ir}}

\vspace{0.4em}
\noindent\textbf{علی جباری}\\
گروهِ مهندسیِ کامپیوتر، دانشگاه ارومیه

\end{document}
